\chapter*{\large ЗАКЛЮЧЕНИЕ}
\addcontentsline{toc}{chapter}{ЗАКЛЮЧЕНИЕ}

В главе~\ref{ch:theory} описан сценарий использования модели в связке с программой Lyrata, объяснено, почему игры на Ren'Py~\cite{srcrenpydocs} дают воспроизводимый поток пар ``текст---музыка'', и как текстовые эмбеддинги RuBERT~\cite{kuratov2019,devlin2018} и вектор аудиопризнаков~\cite{librosa} задают числовое представление данных; приведён обзор архитектур от линейной модели до оконного трансформера~\cite{vaswani2017,nikolenko2018}. Глава~\ref{ch:dataset} фиксирует шаги создания датасета: разбор сценария, учёт шумов разметки, связывание с таблицей музыкальных признаков и каталогом кандидатов для проверки подстановки трека. Глава~\ref{ch:training} задаёт постановку обучения по средней квадратичной ошибке, определения метрик на тесте и метрики попадания в топ доли кандидатов; там же описан протокол \texttt{leave\_games\_out}~\appa{5} и результаты сравнения сохранённых прогонов.

Основной практический вывод: все рассмотренные нейросетевые варианты дают долю попаданий заметно выше случайного уровня при $p = 35$\%, но далеки от насыщения; лучшие значения cosine-hit порядка $0{,}61$ при этом означают, что автоматический список рекомендаций будет часто ошибаться без возможности правки пользователем.

Разрыв между качеством на валидации и на новых играх показывает, что для решений по приложению нужно опираться прежде всего на тест с отложенными играми. Усложнение архитектуры само по себе не даёт гарантированного выигрыша.

Для Lyrata модель имеет смысл как вспомогательный инструмент ранжирования или фильтрации, а не как полная замена выбора музыки~\cite{nikolenko2018}. Перспективы улучшения связаны с объёмом и качеством разметки, постановкой обучения ближе к задаче ранжирования и дополнительными проверками восприятия звука при чтении.
